% GENERATED FILE. Edit canonical lessons, then run npm run generate:books.
% canonical-source-hash: fnv1a64:ffb8eaf78b909db7
% canonical-lessons: ES-C67-repaso-a1, ES-C67-repaso-lo-que-reemplazo-los-casos, ES-C67-primera-lectura, ES-C67-segunda-lectura, ES-C67-tercera-lectura

\chapter{Review --- Four Last Pieces, and the Thread}
\label{ch:repaso-a1}
\hlchaptermodality{266}

\begin{chapteropening}
\textbf{By the end of this chapter:} \emph{I can use the last of the A1 structures, and name the thread that runs through them.}

I can use the last of the A1 structures, and name the thread that runs through them.
\end{chapteropening}

\section[(review)]{(review)}

\label{lesson:ES-C67-repaso-a1}



\begin{quote}
The first five in order. (\emph{Primero, segundo, tercero, cuarto, quinto}.)
\end{quote}

\begin{grammarlens}[title={What you've built}]
Four chants.

\textbf{First.} \emph{El primer libro}, \emph{la primera casa} — ordinals go in front, and two of them drop their \emph{o} before a masculine noun.

\textbf{Second.} \emph{Uno es grande, otro es bueno} — a pair split across the sentence. \emph{Otro libro}, with no article in front of it.

\textbf{Third.} \emph{Comer es bueno} — the infinitive naming an activity, where English would reach for \emph{-ing} and Spanish will not.

\textbf{Fourth.} \emph{Hoy como en casa} — the when and the where can move, and the front is where the emphasis lives.
\end{grammarlens}

\begin{grammarlens}[title={Grammar Lens: the thread}]
\emph{Primer} is the \textbf{third} time you have watched Spanish bite the ending off a word leaning on another — after \emph{muy} and \emph{mal}. Three sightings make a habit, not a list of exceptions, and you will meet a fourth without being surprised.

And the moving pieces close a thread that has run through this whole level. Latin marked who did what with \textbf{endings}. Those endings fell away, and Spanish built replacements out of small audible things: an article to say which one, a preposition to mark the person on the receiving end, a comma to address somebody, and a fixed front position to carry emphasis.

Every one of those was a chapter of its own. Together they are one story, and you now have all of it.
\end{grammarlens}

\subsection*{Your turn}
\begin{itemize}
\raggedright
  \item \emph{Say it:} the four chants, in order
\end{itemize}

\emph{Twice through:}

\begin{tcolorbox}[breakable,colback=teal!4,colframe=teal!35!black,title={Before you move on}]
Where does emphasis live in a Spanish sentence? (At the \textbf{front}.)

And what one event explains the article, the personal \emph{a}, the vocative comma and that emphatic position? (Latin's \textbf{endings falling away}.)
\end{tcolorbox}
\section[(review)]{Review — one name, four different jobs}

\label{lesson:ES-C67-repaso-lo-que-reemplazo-los-casos}



\begin{quote}
\textquotedblleft{}I see María.\textquotedblright{} (\emph{Veo a María}.) \textquotedblleft{}María, come.\textquotedblright{} (\emph{María, ven}.)
\end{quote}

\begin{grammarlens}[title={What you've built}]
The name does not change. \emph{María} is \emph{María} in both. Everything that distinguishes them sits \textbf{outside} the word.

\begin{quote}
\emph{Veo} \emph{\textbf{a}} \emph{María.} — she is what the seeing lands on. \emph{María\emph{\textbf{\emph{,}} }ven.} — she is the one being spoken to.
\end{quote}

A preposition in the first, a comma in the second. Latin did both with endings on \emph{María} herself, and had a whole case for each.

The article does a third job: a name points at one person already, so \emph{el} has nothing left to do — \emph{María}, not \emph{la María}. Give her a \textbf{title} and it comes back: \emph{la señora García}.
\end{grammarlens}

\begin{culture}
The fourth job went to \textbf{position}, which decides form as well as meaning.

\begin{quote}
\emph{Mucha agua.} — leaning on a noun, so it agrees. \emph{Come mucho.} — leaning on a verb, so it freezes.
\end{quote}

Same word, two shapes, and the only difference is what it stands next to. The ordinal does this one slot over: \emph{el primer libro}, in front, shortened.

And because the verb's ending already says who did what, the pieces that say \textbf{when} and \textbf{where} can move:

\begin{quote}
\emph{Hoy como en casa.}  ·  \emph{Como en casa hoy.}  ·  \emph{En casa como hoy.}
\end{quote}

Three placements, one meaning — and the front slot is the emphatic one. Latin bought that freedom with case endings; Spanish lost them and bought it back with word order. Same trade in \emph{¡\textbf{Qué} grande!}: the force sits at the front.
\end{culture}

\subsection*{Your turn}
\begin{itemize}
\raggedright
  \item \emph{Say it:} \emph{hoy como en casa}, \emph{como en casa hoy}, \emph{en casa como hoy}
\end{itemize}

\begin{tcolorbox}[breakable,colback=teal!4,colframe=teal!35!black,title={Before you move on}]
What replaced Latin's ending for the person on the receiving end? (A little \emph{\textbf{a}}.) Which slot carries emphasis? (The \textbf{front} one.) And why does \emph{mucho} refuse to agree in \emph{come mucho}? (It leans on a \textbf{verb}.)
\end{tcolorbox}
\section[(reading)]{(reading) — the first paragraph}

\label{lesson:ES-C67-primera-lectura}



\begin{quote}
Everything you are about to read, you already know. Not one word in it is new.

What \textbf{is} new is that there are forty-seven of them in a row.
\end{quote}

\begin{tcolorbox}[breakable,skin=enhanced,colback=orange!6,colframe=orange!45!black,boxrule=0.5pt,arc=1mm,left=6pt,right=6pt,top=4pt,bottom=4pt,fonttitle=\bfseries,title={Read this}]
\begin{quote}
María vive en una casa grande con su hermano. Hoy no trabaja, porque es domingo. Por la mañana come pan. Luego va a leer un libro. El libro es muy bueno, pero es bastante grande. Ahora está cansada. Va a dormir. Mañana habla español con su estudiante.
\end{quote}

Read it once without stopping. Do not translate as you go — let the sentences arrive.

Now read it again, and notice how little work it took.
\end{tcolorbox}

\begin{culture}
Look at what is holding it together, because it is not vocabulary.

\emph{\textbf{Porque}} joins a fact to its reason. \emph{\textbf{Pero}} holds two true things apart. \emph{\textbf{Luego}} puts one thing after another. \emph{\textbf{Va a}} reaches forward into tomorrow. Each of those was a chapter, and each was small.

Eight sentences across six lines is not much. But the longest run of Spanish anywhere in this book so far was a single line of about fifteen words, and nearly every one of those was a drill — a list of verb forms, or one phrase set against another — rather than something a person would actually say.

This is the first time the book has asked you to hold a whole paragraph in your head at once — and the reason it can is that everything in it was paid for, one chapter at a time.
\end{culture}

\subsection*{Your turn}
\begin{itemize}
\raggedright
  \item \emph{Say it:} the paragraph aloud, once, without stopping
\end{itemize}

\emph{Twice through:}

\begin{tcolorbox}[breakable,colback=teal!4,colframe=teal!35!black,title={Before you move on}]
How many words in that paragraph were new to you? (\textbf{None.})

And what made it readable? (Everything in it had already been \textbf{paid for}, one chapter at a time.)
\end{tcolorbox}
\section[(reading)]{(reading) — the second paragraph}

\label{lesson:ES-C67-segunda-lectura}



\begin{quote}
The last paragraph was about María. This one is about her brother, on a different day.

Same rule as before: not one word in it is new.
\end{quote}

\begin{tcolorbox}[breakable,skin=enhanced,colback=orange!6,colframe=orange!45!black,boxrule=0.5pt,arc=1mm,left=6pt,right=6pt,top=4pt,bottom=4pt,fonttitle=\bfseries,title={Read this}]
\begin{quote}
El lunes María trabaja. Su hermano no. Está en casa con el perro. Por la tarde bebe café y escribe a su madre. No hay pan, pero hay agua. Luego va a salir, porque quiere ver a su padre. Por la noche los dos comen en casa.
\end{quote}

Once through, without stopping.

Again — and this time notice that you did not have to work out who \emph{Está} refers to. Spanish left the name out and you kept it anyway.
\end{tcolorbox}

\begin{culture}
\textbf{Su hermano no.} Two words, and the verb is gone. Spanish drops what the sentence before it already said, and the reader carries it forward.

That is the thing a paragraph can do that a sentence cannot. A sentence has to say everything. A paragraph gets to lean on the one before it.

It is also why reading gets \emph{easier} as it gets longer, which is the opposite of what most people expect.
\end{culture}

\subsection*{Your turn}
\begin{itemize}
\raggedright
  \item \emph{Say it:} the paragraph aloud, once, without stopping
\end{itemize}

\emph{Twice through:}

\begin{tcolorbox}[breakable,colback=teal!4,colframe=teal!35!black,title={Before you move on}]
In \emph{Su hermano no}, what word is missing? (\textbf{Trabaja} — and you supplied it without being asked.)

How many new words were in the paragraph? (\textbf{None}, again.)
\end{tcolorbox}
\section[(reading)]{(reading) — the third paragraph}

\label{lesson:ES-C67-tercera-lectura}



\begin{quote}
This one is longer than both of the others. You are not going to notice.
\end{quote}

\begin{tcolorbox}[breakable,skin=enhanced,colback=orange!6,colframe=orange!45!black,boxrule=0.5pt,arc=1mm,left=6pt,right=6pt,top=4pt,bottom=4pt,fonttitle=\bfseries,title={Read this}]
\begin{quote}
El sábado María y su hermana van a casa de su madre. Su madre es profesora y habla español muy bien. Por la mañana comen pan y beben café. Luego María lee un libro y su hermana escribe. Por la tarde caminan hasta la casa de su hermano. Ahora es tarde, pero están bien. Mañana es domingo, y María no trabaja.
\end{quote}

Once through, without stopping.

Now answer this before you read it again: how many people are in it?
\end{tcolorbox}

\begin{culture}
Four people, and Spanish named each one once. After that they are \emph{su madre}, \emph{su hermana}, \emph{su hermano} — held in place by a single small word you learned a long time ago.

Count what you actually did here. Sixty-one words. Four people. Two days. One page of grammar you already had.

The first paragraph in this chapter was forty-seven words and felt like a wall. This one is longer and did not. Nothing was added in between except two more paragraphs — which is the only thing that ever fixes this.
\end{culture}

\subsection*{Your turn}
\begin{itemize}
\raggedright
  \item \emph{Say it:} the paragraph aloud, once, without stopping
\end{itemize}

\emph{Twice through:}

\begin{tcolorbox}[breakable,colback=teal!4,colframe=teal!35!black,title={Before you move on}]
Which of the three paragraphs was hardest? (\textbf{The first} — and it was the shortest.)

What changed? (\textbf{Nothing was added.} You simply read three of them.)
\end{tcolorbox}
